\noindent\textbf{摘要}\quad
银行提额后，客户可能欠款更多，使用率却更低。贷后人员既要识别新增风险，也要辨清评分变化的来源，并在有限名额内持续关注高风险客户。基于384万条账户月记录，本项目提出\textbf{参考尺度、近稳更新与容量保护协同的信用评分方法}：按固定参考额度重算历史特征，持续接收最新账务；对近稳客户控制评分波动，再以必要修正保留基础模型在预定容量内的入选成员。同一评分由此承接风险排序和跨期比较。提额月2,907名客户的AUC为0.8171；两期559个复核名额覆盖61名未来三个月出现早期逾期的独立客户，较原限幅更新多覆盖2名事件客户。在还款状态或支付结构相近的客群中，最极端5\%评分变化的平均幅度较强风险集成降低39.5\%--68.8\%。六页面工作台进一步连接名单分派、账务核对、证据Agent与案件接续，使评分结果转化为可办理、可交接的贷后任务。

\par\noindent\textbf{关键词：}信用风险；提额管理；稳定评分；容量保护；智能复核

\section{业务背景与项目目标}
信用卡提额以后，贷后人员需要判断：客户的偿付压力是否增加，评分变化是否足以调整复核优先级。例如，额度从1万元提高到1.3万元、余额从6,000元增至7,000元，客户多欠了1,000元，使用率却由60\%降至53.8\%。这一读数同时包含余额变化和额度扩张的影响。提额后借贷行为\textsuperscript{[1]}、决策反馈与数据分布变化\textsuperscript{[2]}的研究，也揭示了银行动作与客户行为之间的联系。

难点随之进入日常作业：策略岗需要解释评分为何改善，复核岗需要决定先查谁，跨月再次入选的客户还需要沿用已核实的材料。长期冻结旧分会压住新风险，删除额度会损失横截面信息；直接限制跳分，则可能改变名单边界，让原本应当关注的客户落选。\textbf{本届赛题同时要求风险区分度与跨期评分一致性\textsuperscript{[15]}。本项目将“稳定更新后仍保留既定风险名单”纳入评分规则，把识别能力、评分可比性和有限复核资源连接起来。}

\subsection{方法创新与应用定位}
方法设计着眼于两个相连的问题。\textbf{评分如何可比：}在一个参考周期内固定额度，从原始历史重算全部相关特征，保留最新余额、支用与还款信息。\textbf{稳定如何兼顾风险关注：}先按账务变化形成近稳候选分，再保护预定名额内的基础模型入选成员；二者冲突时，作最小必要放宽。最终评分统一接受区分度、概率损失、风险覆盖与尾部波动评价。

应用设计进一步回答\textbf{分析如何进入办理}。以客户案件串联名单、账务、动作凭证与处理意见，经办人员沿评分找到证据，完成核对后交给下一责任岗；客户跨月再次入选时，新账务与原有记录一并呈现。项目由此形成从评分更新到持续复核的业务方案，服务额度经营频繁、复核名额有限的贷后团队。

\clearpage
\section{数据获取与业务发现}
\subsection{从384万条账务重建提额前后的客户状态}
主数据来自公开信用卡账户月记录，共3,840,312行、103,558名客户\textsuperscript{[3]}。以观察到的正额度增加识别提额事件，仅用评分时及以前账务生成特征，按客户隔离拟合、校准与评价角色。提额月纳入2,907人，相邻月纳入2,677人，评分时均未逾期。以未来三个月已观察DPD$>0$为早期风险终点，两期分别有2,841、2,623人结果可判，观察到99、88例事件。未知结局保留为空，客户仍占用复核名额；区分度和概率损失在已知结局上计算。

\begin{table}[H]\centering\caption{数据来源及其支持的业务分析}
\fontsize{9.1bp}{12.5bp}\selectfont
\begin{tabularx}{\linewidth}{@{}p{24mm}p{36mm}YY@{}}
\toprule
来源 & 原始规模 & 时间与终点 & 研究用途 \\
\midrule
主信用卡面板 & 103,558客 / 384万行 & 相对月份；三个月DPD>0 & 提额后的风险与稳定性 \\
独立额度面板 & 38,531客 / 14个月 & 2024.03–2025.04；无违约标签 & 余额与使用率变化 \\
拍拍贷LC/LP & 328,553标 / 320万期 & 实际应还与实还日期 & 按日历时间检验还款风险 \\
Prosper & 113,937原始记录 & 前12还款周期核销 & 核销风险识别 \\
新颜作者基准 & 85,236客 / 42特征 & 作者还款异常标签 & 还款异常模型比较 \\
\bottomrule\end{tabularx}
\tabnote{各来源按真实终点独立评价，时间窗口与切分规则见正文及附录。}
\end{table}


\subsection{余额增加时，使用率仍可能下降}
独立额度面板覆盖38,531名客户、14个月\textsuperscript{[4]}。在26,816次额度与余额同时增加的转移中，9,710次使用率反而下降，占36.21\%，客户簇95\%区间为[35.60\%,36.81\%]。留出客户和每客首次提额两种口径下，该比例分别为35.59\%、43.17\%。\textbf{同样的使用率改善，可能对应不同的欠款变化；贷后判断需要同时读取余额和额度。}

\fig{v17_external_mechanism}{提额如何改变使用率读数。a为余额增加而使用率下降的比例及客户簇95\%区间；b将同一新余额分别除以当期额度和提额前额度，比较两种尺度。}

全部75,597次提额转移中，13,207次新余额超过旧额度，占17.47\%。据此，工作台并列呈现真实使用率与参考读数：实际超限按当前额度判断，参考比率用于跨期比较。主研究按客户相对月份跟踪账户演化；其他来源按各自日期和真实终点独立评价。机构接入时，再从审批与策略凭证中确认调整原因、发起方及新增资质信息，为参考周期管理提供依据。

\clearpage
\section{评分方案：兼顾跨期可比与风险名单}
\subsection{固定参考周期，持续接收最新账务}
令$H_t$为评分时可见的非额度历史，$l_e>0$为周期$e$的参考额度。特征构建器从原始历史按$l_e$重算全部额度依赖字段，余额、支用、还款及缺失状态随账务更新。244维XGBoost与414维LightGBM等权融合，经独立客户校准得到基础预测$p_t$；最终评分$q_t$同时用于风险排序与跨期比较，$Q_t=100(1-q_t)$为百分制显示。

\fig{v18_framework}{稳定评分与客户复核的总体架构。最新账务经参考尺度重算，形成基础预测；近稳更新与容量保护共同生成最终评分$q_t$。真实额度、敞口及动作凭证支持证据核对与案件办理。}

\subsection{先形成近稳候选，再保护既定名额}
近稳条件$G_t$要求同账户、同模型、同参考周期，相邻两期未逾期且历史充足；余额、总支付、总支用、最低应还及三项应收金额完整，变化均不超过参考额度的5\%。取$r=0.01$，将基础预测投影至前次评分$q_i^{-}$附近，得到候选分
\[
 z_i=\begin{cases}\Pi_{[\max(0,q_i^{-}-r),\,\min(1,q_i^{-}+r)]}(p_i),&G_i=1,\\p_i,&G_i=0.\end{cases}
\]
在完整批次中预定名额$k$，令$A=\operatorname{Top}_k(p)$，$\tau$为基础预测第$k$与第$k+1$名之间的分界。将入选者的候选分投影到$[\tau^{+},1]$，其余投影到$[0,\tau]$，得到最终$q_i$，满足
\[
\operatorname{Top}_k(q)=\operatorname{Top}_k(p).
\]
候选分符合名单边界时保持原值，发生冲突时移至最近边界。该规则在给定分界区间内作最小平方修正，保留既定容量下的入选成员。主版本保护10\%容量，六容量版本支持分级复核；候选客户或名额改变时整批重算。

\subsection{按动作证据管理参考周期}
经核实的经营性提额沿用原周期，资质重审、账户变更或模型换版触发周期重建，原因未知时先补充凭证。固定客户历史、前次状态、批次与名额后，纯额度改变不影响参考评分；新的账务则持续驱动更新。附录给出并列处理、最小修正证明和批次接口。

\clearpage
\section{风险识别：与强模型在同一批客户上比较}
\subsection{同一评分、同一名额，检验风险覆盖}
提额月AUC为0.8171，291个名额覆盖99例已观察早期事件中的46例，名单命中率15.8\%；相邻月AUC为0.7701，268个名额覆盖88例中的33例。对照包括开发期按AP选出的风险集成S、尾部优化T、目标加权R、参考混合FX与参考集成E；同底模静态和原1pp更新用于分离更新机制的贡献。

\begin{table}[H]\centering\caption{10\%复核容量下的风险识别与事件覆盖}
\fontsize{9bp}{12bp}\selectfont
\begin{tabularx}{\linewidth}{@{}Yrrrrrr@{}}\toprule
 & \multicolumn{3}{c}{提额月：2,907人 / 99事件} & \multicolumn{3}{c}{相邻月：2,677人 / 88事件}\\方法 & AUC & AP & 覆盖 & AUC & AP & 覆盖\\\midrule
目标加权 R & 0.8123 & 0.1455 & 48 & 0.7676 & 0.1215 & 32\\
风险集成 S & 0.8193 & 0.1524 & 44 & 0.7746 & 0.1324 & 34\\
尾部优化 T & 0.8164 & 0.1544 & 43 & 0.7714 & 0.1399 & 32\\
参考混合 FX & 0.8171 & 0.1501 & 45 & 0.7719 & 0.1466 & 32\\
参考集成 E & 0.8188 & 0.1498 & 44 & 0.7737 & 0.1348 & 33\\
同底模静态 & 0.8174 & 0.1500 & 46 & 0.7702 & 0.1435 & 33\\
原1pp更新 & 0.8168 & 0.1492 & 44 & 0.7700 & 0.1454 & 32\\
\textbf{本文方法} & \textbf{0.8171} & \textbf{0.1503} & \textbf{46} & \textbf{0.7701} & \textbf{0.1469} & \textbf{33}\\
\bottomrule\end{tabularx}
\tabnote{名额分别为291、268人；未知结局占用名额。粗体标识本文方法。同底模静态、原1pp更新与本文共用特征、融合及校准。}
\end{table}

\fig{v18_primary_capacity}{有限复核名额下的风险覆盖。内圈以统一0--99刻度展示提额月事件数，外圈为名单命中率；右侧展开10\%容量。两期六档容量的完整结果见附录A.4。}

\subsection{稳定更新后，保留应当进入复核的成员}
同底模静态在两期分别覆盖46、33例事件，原1pp更新降至44、32例。容量保护恢复这3个事件窗口，名额仍为291与268。它同时处理两类名单变化：客户自身风险上调受到限幅，以及其他客户分数变化造成的挤出。\textbf{在预定容量下，最终评分保留基础预测的全部入选成员，将平稳更新与风险关注落实到同一份名单。}

各方法的优势对应不同评价目标：S两期AUC为0.8193、0.7746，R两期覆盖80个事件窗口；本文覆盖79个窗口，并进一步考察尾部波动、独立事件客户与跨月接续。这些结果来自同一历史队列上的冻结模型比较，客户配对区间及概率校准结果见附录。

\clearpage
\section{跨期稳定性：从账务近稳扩展到业务状态相近}
\subsection{用还款状态与交易结构刻画真实月度变化}
在四字段金额近稳组之外，设置两套评分时可观察的业务分组。\textbf{还款与交易状态组}要求支付覆盖状态、交易笔数档位相同；\textbf{支付与交易结构组}要求支付覆盖比、取现份额和笔数相近。它们分别从偿付状态和交易结构衡量客户变化。以$D_t=100|q_t-q_{t-1}|$度量评分波动，CVaR95为幅度最大的5\%记录的平均变化，单位为风险概率百分点。

\begin{table}[H]\centering\caption{业务相近客户的极端评分波动（CVaR95）}
\fontsize{9bp}{12bp}\selectfont
\begin{tabularx}{\linewidth}{@{}Ylrrrr@{}}\toprule
业务口径 & 观察期 & 配对记录 & 本文 & 风险集成 S & 尾部优化 T\\\midrule
账务金额近稳 & 提额月 & 582 & 2.184 & 10.994 & 3.634\\
 & 相邻月 & 439 & 1.506 & 4.299 & 1.535\\
\addlinespace[2bp]
还款与交易状态 & 提额月 & 1063 & 5.288 & 8.735 & 6.003\\
 & 相邻月 & 943 & 4.384 & 7.376 & 5.055\\
\addlinespace[2bp]
支付与交易结构 & 提额月 & 403 & 3.122 & 10.019 & 4.789\\
 & 相邻月 & 284 & 3.173 & 7.253 & 3.780\\
\bottomrule\end{tabularx}
\tabnote{单位：风险概率百分点。各业务组可重叠；完整分组定义、门控重合及配对区间见附录A.9。}
\end{table}

\fig{v18_tail_evidence}{真实月度评分变化的尾部分布与配对差异。a为提额月582对账务近稳记录；b为本文相对五个对手的CVaR95差值，负值表示波动减小。粗线为逐项95\%区间，细线为三客群、两期、七对手共42项比较的同时95\%区间。}

\subsection{保持风险名单，扩大尾部改善的适用客群}
账务近稳组两期P95均为1个百分点，CVaR95为2.184、1.506，较S分别降低80.1\%、65.0\%。还款状态组的降幅为39.5\%、40.6\%，支付结构组为68.8\%、56.2\%；后四项差值的同时95\%区间均低于0。两类业务组与七字段近稳条件的重合率仅为34.5\%--68.0\%，\textbf{说明尾部改善延伸到了按还款状态和支付结构识别的客群}。

5,584次更新中，14次为保持名单而放宽候选边界，占0.251\%。P95定位变化分布的95\%分位，CVaR95刻画其上尾平均幅度；两者结合逐户放宽幅度，呈现常见变化与极端变化。工作台同步保存基础预测、候选分、最终分及放宽原因，使每次调整都能回到账务核对。

\clearpage
\section{连续复核：风险覆盖与客户办理一并接续}
\subsection{两期559个名额覆盖61名独立事件客户}
两期按10\%容量分别安排291、268个名额，合计涉及331名独立客户，覆盖61名已观察早期逾期客户。228人连续入选，相邻月新增40人。与尾部优化T相比，多覆盖6名事件客户、涉及客户少14人；与目标加权R相比，覆盖相同的61名事件客户、涉及客户少25人。再次入选时，已核实的身份、动作材料和办理意见随原案接续，经办人员集中核查新增账务。

\fig{v18_continuity_capacity}{连续两期的客户覆盖与复核范围。a为六档容量下覆盖的独立事件客户数；b比较同样559个名额涉及的独立客户与事件客户。再次入选时核查新账务，并沿用此前已确认的材料。}

\subsection{受限更新中的风险识别与概率校准}
以原方案真正受限的106次更新固定评价子组，其中105次结果可判，含12个事件。容量保护相对原更新使AUC从0.7079升至0.7384，Brier从0.09832降至0.09704。该子组平均预测为8.74\%，实际事件率为11.43\%；8次阳性向上受限记录的最终评分仍低于静态预测，因此将概率校准与名单覆盖分别纳入持续监测。

整批名单共恢复3个事件窗口，包含未直接触发限幅、却被其他客户挤出的记录。14次必要放宽中，最大额外幅度为18.958个百分点：该客户基础风险已明显下降，原限幅使评分仍处于较高位置；容量保护将其移回基础名单的未入选一侧，为原入选成员保留名额。逐行变化、条件区间及6／12个月DPD30+随访见附录A.5。

\subsection{全量客户分流到明确的办理队列}
工作台覆盖最初2,907名客户。其中2,677人进入相邻月完整评分批次，安排268个复核名额；189人保留补齐后的评分记录，40人进入逾期核对，1人接续补账。按办理任务划分为互斥的五类队列：名单复核268人、差异核对6人、持续观察2,592人、逾期核对40人、账务补齐1人。评分批次对应风险筛选，全量队列对应客户管理；候选成员变化后按新批次重新评分。

\clearpage
\section{产品应用：从风险名单到案件接续}
六个主页面围绕同一客户展开。负责人从总览识别重点客群，在风险队列分派任务；经办人员核对账务与评分变化，补齐动作凭证，由AI协同形成带来源的意见，再交给下一责任岗。页面共享客户、观察月和案件版本，\textbf{将“先查谁、查什么、交给谁”组织成连续的复核流程。}

\fig{workbench_current}{六页面复核工作台及案件流转示例。a总览客群与任务分流；b按风险序位分派案件；c核对账务、评分与名单保护；d登记动作凭证和办理意见；e生成带来源的Agent草稿；f跟踪责任岗、跟进日期及案件版本。}

\clearpage
\section{商业落地：嵌入贷后流程，按作业价值验收}
\subsection{以一支贷后团队为起点，连接现有系统与岗位}
成熟客户管理平台将统一客户视图、额度管理与持续复核相连接，作业型Agent进一步提供优先队列、账户上下文和办理建议\textsuperscript{[11,13]}。本方案聚焦\textbf{提额后复核与跨月跟进}：策略岗确定名额与参考周期，复核岗沿评分核查账务，客户运营岗接续待补材料。以账户和案件为主键连接授信、工单及身份系统，沿用机构现有责任分工。

\begin{table}[H]\centering\caption{产品能力、岗位任务与验收指标}\small
\begin{tabularx}{\linewidth}{@{}p{20mm}YY@{}}\toprule
岗位 & 日常工作中的用法 & 验收指标\\\midrule
贷后复核 & 同页查账、核对凭证、续办原案 & 完整正确率、每件合格材料工时\\
策略运营 & 设置名额、解释评分、管理参考周期 & 名单一致性、风险覆盖、尾部波动\\
客户运营 & 接续待补材料与办理意见 & 退回率、重复联系、接续时长\\
技术与模型 & 接入账务、批次评分、工单回执 & 重算一致性、版本管理、异常恢复\\\bottomrule
\end{tabularx}\end{table}

首个落地单元是一支负责提额后跟进的贷后团队。评分适配器接入账务与批次名单，工作台支持取数、查证和办理，案件服务保存版本与回执。先以只读接入开展并行复核，再接通统一身份、岗位权限及工单，按需接入机构语言模型。\textbf{评分、案件和语言服务分层，使模型升级与流程扩展可以独立实施。}商业交付按实施接入与年度组件维护设计：前者完成字段映射和工作流配置，后者支持升级、监测及运维；采购验收对应处理质量、接口可靠性和实际作业成本。

\subsection{用合格材料和全流程工时衡量投入产出}
作业试验已配置18个正式账务案例与3个练习案例。A提供账务表，B增加常规图表，C增加参考尺度、变化分解、来源草稿与接续事项，三组共享底层证据。每人办理9案、每种界面3案，同人不重复同案；按岗位与经验平衡顺序，记录提交、盲评、退回及再次办理。

首轮现场试点拟招募24名从业者，以C对B为主比较，估计配对差异与方差。验收同时考察读数正确率、材料完整正确率和每件合格材料总工时，失败与返工全部计时。按年案件量$N$、实测每案节时$\Delta t$（分钟）、全成本小时费率$w$及年度总成本$C$计算净作业价值$V=N\Delta t\,w/60-C$，用于确定接入与维护预算。新机构客户按预登记预测和成熟结局进行一次性评价，分别验收风险效果与作业效果。

\section{结论}
本项目围绕提额后的评分比较、客户筛选和案件接续形成完整方案。统一参考尺度持续接收新账务，近稳更新与容量保护共同生成最终评分，在保持预定容量内基础风险名单的同时，收缩真实月度变化的尾部波动。两期复核覆盖61名独立早期逾期客户，工作台将账务、动作凭证和处理意见纳入同一案件。\textbf{面向额度经营频繁的贷后团队，本方案提供可比的风险读数、连续的客户记录和有据可查的办理过程。}
\label{bodyend}
